% 摘自 mathexam.dtx 文档中的题型示例（v2.3.1）
\documentclass{article}
\usepackage[nospace]{mathexam}
\course{高等代数}
\university{西南大学}
\school{数学与统计学院}
\begin{document}
\makehead
\begin{makepart}{判断题}[2]
\begin{problem}
  假设直角三角形的两直角边边长分别为3和4, 则斜边长度为5. \true
\end{problem}
\end{makepart}
\begin{makepart}{选择题}[2]
\begin{problem}
  假设直角三角形的两直角边边长分别为3和4, 则斜边长度为\pickout{C}
\begin{abcd}
 \item 7;
 \item 6;
 \item 5;
 \item 4.
\end{abcd}
\end{problem}
\end{makepart}
\begin{makepart}{填空题}
\begin{problem}[2]
  假设直角三角形的两直角边边长分别为3和4, 则斜边长度为\fillin{5}.
\end{problem}
\end{makepart}
\begin{makepart}{计算题}
\begin{problem}[10]
  假设直角三角形的两直角边边长分别为3和4, 则斜边长度为?
\begin{solution}[10em]
  根据勾股定理.\score{5}
  因此斜边为$\sqrt{3^2+4^2}=5$.\score{5}
\end{solution}
\end{problem}
\end{makepart}
\end{document}
